\chapter{Conclusão}\label{cap:conclusao}
\section{Resposta aos objetivos}
O trabalho resultou na construção e documentação de um protótipo modular que integra mecanismo de disco, ESP32, Raspberry Pi, comunicação I²C, serviços de aplicação e interface de cuidado. Os materiais recebidos permitem demonstrar montagem física, operações registradas entre central e placas, atualização de software e evidências parciais de transporte de unidades. O capítulo de visão passa a explicitar aquisição, processamento e decisão temporal, mantendo a divergência de implantação em aberto. \cite{DC_VALIDACAO2026,DC_VISAO2026}

Esses resultados não encerram o objetivo de dispensação autônoma confiável. Os registros vinculados acrescentam observações humanas parciais, mas não uma campanha independente que estime a taxa pedido = saída; houve retenção em ciclos registrados, o RTC não está ativo e a agenda automática permanece desligada. Também não se dispõe de avaliação rotulada da visão, campanha completa do agente ou avaliação de uso que sustente melhoria de adesão ou eficácia clínica. \cite{DC_VALIDACAO2026,DC_VISAO2026}

\section{Contribuições e continuidade}
As contribuições demonstradas concentram-se na integração tecnológica, na rastreabilidade das operações e na identificação de condições que exigem intervenção humana. A distinção entre comando reconhecido, sessão encerrada, leitura da cavidade e quantidade coletada é necessária para que o acompanhamento não transforme incerteza em sucesso aparente.

O fechamento do trabalho depende da campanha física com referência independente, da caracterização elétrica sob carga, do comissionamento temporal e da reconciliação de versões das frentes de visão e agente. A utilização pretendida com pessoas requer ainda avaliação institucional, privacidade e critérios de segurança compatíveis. O protótipo é uma plataforma acadêmica em desenvolvimento, não um produto clinicamente validado.
\pendencia{CON01}{Revisar a resposta aos objetivos após a campanha final, substituindo apenas os limites efetivamente resolvidos e conservando a rastreabilidade das evidências.}
\pendencia{CON02}{Validar com a equipe as contribuições e a continuidade, sem apresentar funcionamento autônomo, conformidade ou benefícios clínicos que não foram demonstrados.}
